\documentclass{article}
\usepackage{amsmath}
\begin{document}
\section*{Local timers with shared clock support on the same cave history}
Both observers' old calibration blocks the first audit window;
their replacement blocks the second. Both replacement duties
must act within a fixed local service interval on the modeled
mission clock. Merely issuing identical numeric instructions does
not establish either an on-time switch or a common side of
the outstanding audit window.

\paragraph{Source contract.}
At the actual qualified cave rest position, the service cuts
are ticks 2 and 4, with a tick of 15.625 ms. The requested
effect is bounded to ticks 1 through 7. Holds, native arc
qualification and local gate evidence are checked in the same
older cave replay. Two distinct native addressed instructions
arm the local observers at rest. Each actual physical control
callback yields its numeric threshold to its own timer. Each
timer evaluates only its own current clock read. The controlling
adapter observes whether the current stable source actually
binds both timers to the same commissioned clock object, but
cannot read the hidden true offset for use in admission.
This opaque equality comes from the modeled local clock
authority, not equal sender names or jointly reading numeric
values. Distinct commissioned clocks have the same individual
offset alphabet and no asserted coupling. The replay declares
continued availability of these committed clocks through
the whole bounded phase and validates that both arm messages
actually arrive.

\paragraph{Joint continuation.}
The inherited exact support calculation finds a unique common
target 3 when both clocks have one shared offset in
\(\{1,-2\}\). The new schedules then both act at physical tick 2,
or both act at tick 5. The old audit gets the first cut in
one branch or the second in the other. Under the product
support no common target meets both the old audited service
and the earlier/latest new-action bounds. A conservative
controller withholds that numeric action; it still delivers
the old audit but does \emph{not} claim the new duties achieved
their requested window. This is an exact boundary for the
common-target interface, not a prohibition on a stronger
qualified cross-clock intervention. Both arms reserve the
actual native transmit room for the five-byte primary attempt.
Primary begins at tick 7 in every arm, before the same checked
backup and route continuations. This advances rest by 109.375 ms
relative to the parent before forwarding the primary; that
cost is charged to every arm.

\paragraph{Bounded validation.}
Both mapped route geometries, both before/after crash orders in
each of the gate and dock transactions, and all six qualified or
unqualified clock-source worlds are crossed in 48 histories
per arm. The direct finite checker consumes the same observable
scope and ties the interval calculation. Both lawful arms
retain all 48 old audits, establish both timed actions in the
16 histories with a qualified common clock, and preserve every
exact dock result with maximum egress 453.957198 s. Trusting
the old joint qualifier even after its scope no longer holds
runs both timers in all 48, but loses 16 old audits. The
unqualified cases with coincident actual offset are kept:
coincidence in one history does not license an unproved
prospective binding. The underlying pinned middleware
does not provide an external guarantee of a common
clock or atomic remote arming. These are held-in model results,
not independent deployment or a field synchronization claim.
\end{document}
